
\section{Regularisation of Feynman integrals}

\SourcePage{548}{553}

The physical renormalisation conditions considered in~\S\,110 allow, in principle, obtaining unambiguously finite values for the amplitude of any electrodynamic process when computed in any approximation of perturbation theory.

Let us first clarify the character of the divergences in Feynman integrals. Consider an $n$-th order diagram ($n$~vertices) with $N_e$~external electron lines and $N_\gamma$~external photon lines.

The number~$N_e$ is always even; the electron lines in the diagram form $N_e/2$ continuous sequences, each beginning and ending at an external end. The number of internal electron lines in each such sequence is one less than the number of vertices lying on it; therefore the total number of internal electron lines in the diagram equals $n - N_e/2$.

Each internal photon line connects two vertices; their total number is $(n - N_\gamma)/2$.

Each internal photon line carries a propagator $D(k) \sim k^{-2}$. Each internal electron line carries a propagator $G(p) \sim p^{-1}$. The total power of 4-momenta in the denominator of the integrand is thus
\[
2\,\frac{n - N_\gamma}{2} + \frac{n - N_e/2}{1} = 2n - \tfrac{1}{2}N_e - N_\gamma.
\]

The number of independent integrations (over $d^4p$ or $d^4k$) equals the total number of internal lines minus the $n - 1$ additional constraints imposed by 4-momentum conservation at the vertices (one conservation law relates external momenta and does not constrain internal ones). Thus the total number of integrations over all components of 4-momenta is
\[
4\!\left[\frac{n - N_\gamma}{2} + n - \frac{N_e}{2} - (n-1)\right] = 2(n - N_e - N_\gamma + 2).
\]

The degree of divergence $r$ of the diagram is the difference between this number and the power of momenta in the denominator:
\begin{equation}
r = 4 - \tfrac{3}{2}\,N_e - N_\gamma.
\label{eq:112.1}
\end{equation}
Remarkably, this number does not depend on the order~$n$ of the diagram.

The condition $r < 0$ for the diagram as a whole is, however, insufficient for convergence of the integral: every internal block of the diagram must likewise have a negative degree of divergence~$r'$. The presence of blocks with $r' > 0$ can increase the degree of divergence of the diagram as a whole.

\SourcePage{549}{554}

Since $N_e$ and $N_\gamma$ are positive integers, only a few pairs of values $(N_e, N_\gamma)$ give $r \geq 0$. Let us enumerate these cases, immediately excluding $N_e = N_\gamma = 0$ (vacuum closed loops, devoid of physical meaning) and $N_e = 0$, $N_\gamma = 1$ (a ``vacuum current,'' to be discarded according to~\S\,103). The remaining cases are:

% [Feynman diagrams (112.2): five primitive divergent diagrams labeled a through д:
% a) N_e=0, N_γ=2 (r=2): photon self-energy (polarisation operator) — electron loop with two external photon lines
% б) N_e=0, N_γ=3 (r=1): three-photon vertex — triangle with three external photon lines
% в) N_e=0, N_γ=4 (r=0): four-photon scattering (light-by-light) — box with four external photon lines
% г) N_e=2, N_γ=0 (r=1): electron self-energy (mass operator) — electron line with self-energy bubble
% д) N_e=2, N_γ=1 (r=0): vertex correction — triangle with two external electron lines and one photon]
\begin{equation}
\label{eq:112.2}
\end{equation}

In the first case the divergence is quadratic; in all the others ($r = 0$ or $r = 1$) it is logarithmic.

Diagram~\textit{d}) is the first correction to the vertex operator. It must satisfy condition~(\ref{eq:110.19}), which we write in the form
\begin{equation}
\bar{u}(p)\,\Lambda^\mu(p,p;\,0)\,u(p) = 0, \qquad p^2 = m^2,
\label{eq:112.3}
\end{equation}
where
\begin{equation}
\Lambda^\mu = \Gamma^\mu - \gamma^\mu.
\label{eq:112.4}
\end{equation}

Denote the Feynman integral written directly from the diagram rules by $\bar{\Lambda}^\mu(p_2, p_1;\,k)$. This integral diverges logarithmically and does not by itself satisfy~(\ref{eq:112.3}). However, a quantity satisfying this condition is obtained by forming the difference
\begin{equation}
\Lambda^\mu(p_2, p_1;\,k) = \bar{\Lambda}^\mu(p_2, p_1;\,k) - \bar{\Lambda}^\mu(p_1, p_1;\,0)\big|_{p_1^2 = m^2}.
\label{eq:112.5}
\end{equation}

The leading divergent part of the integral~$\bar{\Lambda}^\mu$ is obtained by regarding the virtual photon 4-momentum~$f$ as arbitrarily large in the integrand. It has the form\footnote{The full expression of the integral is given in~\S\,117 (see~(117.2)).}
\[
-4\pi i e^2 \int \frac{\gamma^\nu(\gamma f)\gamma^\mu(\gamma f)\gamma_\nu}{f^2 \cdot f^2 \cdot f^2}\,\frac{d^4f}{(2\pi)^4}
\]
and does not depend on the values of the 4-momenta of the external lines. Therefore in the difference~(\ref{eq:112.5}) the divergence cancels and a finite quantity is obtained. This operation of eliminating the divergence by subtraction is called \emph{regularisation} of the integral.

\SourcePage{550}{555}

Note that regularisability by a single subtraction is ensured by the fact that in this case the divergence is only logarithmic---the weakest possible. If the integral contained divergences of various orders, one subtraction might be insufficient.

Having determined the first correction~$\Gamma^\mu$ (the first term of the expansion of~$\Lambda^\mu$), the first correction to the electron propagator (diagram~(\ref{eq:112.2},\textit{c})) can be computed using the Ward identity~(\ref{eq:108.8}), which may also be written as
\begin{equation}
\frac{\partial \mathcal{M}(p)}{\partial p_\mu} = -\Lambda^\mu(p,p;\,0),
\label{eq:112.6}
\end{equation}
introducing the mass operator~$\mathcal{M}$ in place of~$\mathcal{G}$ and $\Lambda^\mu$ in place of~$\Gamma^\mu$. This equation must be integrated with the boundary condition
\begin{equation}
\bar{u}(p)\,\mathcal{M}(p)\,u(p) = 0, \qquad p^2 = m^2,
\label{eq:112.7}
\end{equation}
which follows from~(110.20).

Finally, for the first term in the expansion of the polarisation operator, we use identity~(108.14); after contraction over pairs of indices it gives
\[
\frac{3}{4\pi}\,\frac{\partial^2 \mathcal{P}}{\partial k_\rho\,\partial k^\sigma} = 2\mathcal{G},
\]
relating the scalar functions $\mathcal{P} = \tfrac{1}{3}\mathcal{P}^\mu{}_\mu$ and $\mathcal{G} = \mathcal{G}^\mu{}_{\rho}{}^{\rho\mu}$. Both depend only on the scalar~$k^2$:
\begin{equation}
2k^2\,\mathcal{P}''(k^2) + \mathcal{P}'(k^2) = \tfrac{4}{3}\,\mathcal{P}(k^2),
\label{eq:112.8}
\end{equation}
where primes denote differentiation with respect to~$k^2$. Together with condition~$\mathcal{P}'(0) = 0$, this also gives
\begin{equation}
\mathcal{P}(0) = 0.
\label{eq:112.9}
\end{equation}

\SourcePage{551}{556}

In the first approximation, $\mathcal{P}(k^2)$ is determined by diagram~(\ref{eq:112.2},\textit{e}) (with 4-momenta of the ends $k$, $k$, $0$, $0$). The corresponding Feynman integral (denote it~$\bar{\mathcal{G}}(k^2)$) diverges logarithmically, and its regularisation is accomplished by a single subtraction using condition~(\ref{eq:112.9}):
\[
\mathcal{G}(k^2) = \bar{\mathcal{G}}(k^2) - \bar{\mathcal{G}}(0).
\]
Then $\mathcal{P}(k^2)$ is determined as the solution of equation~(\ref{eq:112.8}) with boundary conditions $\mathcal{P}(0) = 0$, $\mathcal{P}'(0) = 0$.

In the next approximation, the correction to the vertex operator ($\Lambda^{(2)}_\mu$) is determined by diagrams~(\ref{eq:106.10},\textit{d}--\textit{i}). The irreducible diagram~(\ref{eq:106.10},\textit{d}) is computed by the same regularisation~(\ref{eq:112.5}) as in the computation of the first-approximation~$\Lambda^{(1)}_\mu$. In the reducible diagrams, the internal self-energy and vertex parts of lower order are replaced by the already-known (regularised) first-approximation quantities ($\mathcal{P}^{(1)}$, $\mathcal{M}^{(1)}$, $\Lambda^{(1)}_\mu$), after which the integrals are once again regularised using~(\ref{eq:112.5}).\footnote{In diagrams of still higher approximations it may also be necessary to replace ``four-tails'' by the regularised values of~$\mathcal{G}$.}

The systematic procedure described gives, in principle, finite expressions for $\mathcal{P}$, $\mathcal{M}$, and~$\Lambda_\mu$ in any approximation of perturbation theory. Thereby it becomes possible to compute the amplitudes of physical scattering processes, whose diagrams contain the blocks~$\mathcal{P}$, $\mathcal{M}$, $\Lambda_\mu$ as constituent parts.

We see that the physical conditions established in~\S\S\,110 and~111 turn out to be sufficient for the unambiguous regularisation of all Feynman integrals. This is by no means a trivial property of quantum electrodynamics, and is called \emph{renormalisability}.\footnote{Another approach to the theory of renormalisation in quantum electrodynamics is given in: N.\,N.~Bogolyubov, D.\,V.~Shirkov, \textit{Introduction to the Theory of Quantised Fields}, Moscow: Nauka, 1984.}

\SourcePage{552}{557}

For practical computation of radiative corrections the described procedure may not be the simplest or most rational. In the next chapter we shall see that it is often more convenient to begin by computing the imaginary parts of the corresponding quantities; these are given by integrals that do not diverge and contain no divergent integrals. The full quantity is then determined by analytic continuation using dispersion relations. This avoids the need for direct regularisation by subtraction.
